% !TEX TS-program = lualatex
% https://github.com/Sorbus413/tusurlatex
% https://gitverse.ru/Sorbus413/tusurlatex
% https://gitflic.ru/Sorbus413/tusurlatex

\documentclass[a4paper,14pt]{extarticle}

% Преамбула
\usepackage{fontspec}
\usepackage{pdfpages}
\usepackage[normalem]{ulem}

% шрифты
\defaultfontfeatures{Ligatures=TeX}
\setmainfont{Liberation Serif}
\setmonofont{Liberation Mono}[Scale=MatchLowercase]
\setsansfont{Liberation Sans}
\newfontfamily\cyrillicfont{Liberation Serif}[Script=Cyrillic]
\newfontfamily\cyrillicfontsans{Liberation Sans}[Script=Cyrillic]
\newfontfamily\cyrillicfonttt{Liberation Mono}[Script=Cyrillic,Scale=MatchLowercase]

% язык
\usepackage{polyglossia}
\setdefaultlanguage{russian}
\setotherlanguage{english}

% абзацы и списки
\usepackage{enumerate}
\usepackage{indentfirst}
\setlength{\parindent}{12.5mm}
\usepackage{float}
\usepackage{setspace}
\onehalfspacing

% таблицы
\usepackage{array,tabularx,tabulary,booktabs}
\usepackage{longtable}
\usepackage{multirow}

% tikz
\usepackage{tikz}

% Команды для crow's foot нотации
\newcommand{\crowone}{%
    \tikz[baseline=-0.5ex, x=1.2ex, y=1.2ex]{
        \draw[thick] (0,0) -- (2,0);
        \draw[thick] (1,-.8) -- (1,.8);
        \draw[thick] (1.5,-.8) -- (1.5,.8);
    }%
}

\newcommand{\crowoneorzero}{%
    \tikz[baseline=-0.5ex, x=1.2ex, y=1.2ex]{
        \draw[thick] (0,0) -- (1,0);
        \draw[thick] (0.5,-0.8) -- (0.5,0.8);
        \draw[thick] (1.5,0) circle (0.5);
        \draw[thick] (2,0) -- (2.5,0);
    }%
}

\newcommand{\crowmany}{%
    \tikz[baseline=-0.5ex, x=1.2ex, y=1.2ex]{
        \draw[thick] (0,0) -- (2.5,0);
        \draw[thick] (1.5,0) -- (2.5,-.8);
        \draw[thick] (1.5,0) -- (2.5,.8);
    }%
}

\newcommand{\crowoneormany}{%
    \tikz[baseline=-0.5ex, x=1.2ex, y=1.2ex]{
        \draw[thick] (0,0) -- (2.5,0);
        \draw[thick] (1,-.8) -- (1,.8);
        \draw[thick] (1.5,0) -- (2.5,-.8);
        \draw[thick] (1.5,0) -- (2.5,.8);
    }%
}

\newcommand{\crowzeroormany}{%
    \tikz[baseline=-0.5ex, x=1.2ex, y=1.2ex]{
        \draw[thick] (0,0) -- (0.5,0);
        \draw[thick] (1,0) circle (0.5);
        \draw[thick] (1.5,0) -- (2.5,0);
        \draw[thick] (1.5,0) -- (2.5,-.8);
        \draw[thick] (1.5,0) -- (2.5,.8);
    }%
}

% математика
\usepackage{amsmath}
\usepackage{pdflscape}
\usepackage{amssymb}

% пути изображений
\usepackage{graphicx}
\graphicspath{{./images/}}

% подписи объектов
\usepackage{caption}

% разделитель
\DeclareCaptionLabelSeparator{dash}{\space--\space}

% рисунки
\captionsetup[figure]{justification=centering, labelsep=dash, format=plain} 
\addto\captionsrussian{\renewcommand{\figurename}{Рисунок}\vspace{24pt}}
\renewcommand{\thefigure}{\thesection.\arabic{figure}}

% таблицы
\captionsetup[table]{justification=raggedright, labelsep=dash, format=plain, singlelinecheck=false}
\renewcommand{\thetable}{\thesection.\arabic{table}}

% листинги
\usepackage{listings}

% нумерация в пределах секций
\usepackage{chngcntr}
\counterwithin{figure}{section}
\counterwithin{table}{section}

% гиперссылки
\usepackage{hyperref}
\usepackage{refcount}
\hypersetup{
	colorlinks=true,
	linkcolor=black,
	filecolor=black,      
	urlcolor=black,
	citecolor=black
}

% переносы
\hyphenpenalty=10000
\exhyphenpenalty=10000
\tolerance=9999
\emergencystretch=3em
\hbadness=10000 
\hfuzz=2em

% границы страниц
\usepackage{geometry}
\geometry{left=3cm}
\geometry{right=1.5cm}
\geometry{top=2cm}
\geometry{bottom=2cm}

% списки
\usepackage{enumitem}
\setlist[enumerate,itemize]{leftmargin=0pt, labelindent=15.5mm, labelsep=0.5em, itemindent=0pt, labelwidth=!}
\makeatletter
\AddEnumerateCounter{\asbuk}{\@asbuk}{м)}
\makeatother
\setlist{nolistsep}
\renewcommand{\labelitemi}{--}
\renewcommand{\labelitemii}{--}
\renewcommand{\labelenumi}{\arabic{enumi}.}
\renewcommand{\labelenumii}{\asbuk{enumii})}
\renewcommand{\labelenumiii}{\arabic{enumii})}

% содержание
\usepackage{tocloft}
\renewcommand{\cfttoctitlefont}{\hspace{0.38\textwidth}\bfseries\MakeTextUppercase}
\renewcommand{\cftsecfont}{\hspace{0pt}}
\renewcommand\cftsecleader{\cftdotfill{\cftdotsep}}
\renewcommand\cftsecpagefont{\mdseries}
\setcounter{tocdepth}{2}
\addto\captionsrussian{\renewcommand{\contentsname}{ОГЛАВЛЕНИЕ}}

% нумерация страниц
\usepackage{fancyhdr}
\pagestyle{fancy}
\fancyhf{}
\cfoot{\textrm{\thepage}}
\fancyheadoffset{0mm}
\fancyfootoffset{0mm}
\setlength{\headheight}{17pt}
\renewcommand{\headrulewidth}{0pt}
\renewcommand{\footrulewidth}{0pt}
\fancypagestyle{plain}{
    \fancyhf{}
    \cfoot{\textrm{\thepage}}
}

% заголовки секций в оглавлении
\usepackage{textcase}
\makeatletter
% Флаг для отслеживания, что это настоящая секция
\newif\ifrealsection
\realsectiontrue

% Сохраняем оригинальную команду section
\let\oldsection\section
\renewcommand{\section}{%
    \realsectiontrue
    \oldsection
}

% Переопределяем \addcontentsline для секций
\let\oldaddcontentsline\addcontentsline
\renewcommand{\addcontentsline}[3]{%
    \def\tempa{#1}\def\tempb{toc}%
    \def\tempc{#2}\def\tempd{section}%
    \ifx\tempa\tempb
        \ifx\tempc\tempd
            \ifrealsection
                \oldaddcontentsline{#1}{#2}{\MakeTextUppercase{#3}}%
                \global\realsectiontrue
            \else
                \oldaddcontentsline{#1}{#2}{#3}%
                \global\realsectiontrue
            \fi
        \else
            \oldaddcontentsline{#1}{#2}{#3}%
        \fi
    \else
        \oldaddcontentsline{#1}{#2}{#3}%
    \fi
}
\makeatother

%\renewcommand{\thesection}{\arabic{section}.}  % Точка после номера section в TOC
%\renewcommand{\thesubsection}{\thesection\arabic{subsection}}  % Убираем точку после section в subsection
% заголовки
\usepackage[compact,explicit]{titlesec}
\titleformat{\section}[block]{\centering\bfseries}{}{0pt}{\thesection\quad\MakeTextUppercase{#1}}
\titleformat{\subsection}[block]{\centering\bfseries}{}{0pt}{\thesubsection\quad#1}
\titleformat{\subsubsection}[runin]{}{}{12.5mm}{\thesubsubsection\quad#1.}
\titleformat{\paragraph}[block]{\vspace{24pt}}{}{0pt}{\vspace{12pt}\centering\textbf{#1}}[\vspace{12pt}]

\titlespacing*{\section}{0pt}{0pt}{24pt}
\titlespacing*{\subsection}{0pt}{24pt}{24pt}

\titlespacing*{\section}{0pt}{0pt}{24pt}
\titlespacing*{\subsection}{0pt}{24pt}{24pt}

%\renewcommand{\cftsecpresnum}{ГЛАВА\space} % Слово Глава перед section
\newlength\mylength
\settowidth\mylength{\cftsecpresnum}
\addtolength\cftsecnumwidth{\mylength}

% заголовок 1 с новой страницы
\let\stdsection\section
\renewcommand\section{\clearpage\stdsection}

% секции без номеров
\makeatletter
\newcommand{\anonsection}[1]{
    \newpage
    \phantomsection
    \global\realsectionfalse
    \paragraph{\textbf{\centering{#1}}}
    \addcontentsline{toc}{section}{#1}
}
\makeatother

% приложения
\makeatletter
\newcommand{\appsection}[2][]{%
    \newpage
    \phantomsection   
    \global\realsectionfalse
    \stepcounter{AppCount}
    \setcounter{figure}{0}
    \setcounter{table}{0}
    \setcounter{listing}{0}
    \renewcommand{\thefigure}{\theAppCount.\arabic{figure}}
    \renewcommand{\thetable}{\theAppCount.\arabic{table}}
    \renewcommand{\thelisting}{\theAppCount.\arabic{listing}}
    \begin{center}
        \textbf{Приложение \theAppCount}%
        \ifx&#1&%
        % если аргумент пустой, ничего не добавляем
        \else%
        \\
        (#1)%
        \fi%
        \\
        \textbf{#2}
    \end{center}
    \ifx&#1&%
    \addcontentsline{toc}{section}{Приложение \theAppCount\space #2}%
    \else%
    \addcontentsline{toc}{section}{Приложение \theAppCount\space(#1) #2}%
    \fi%
}
\makeatother

% счётчик приложений
\newcounter{AppCount}
\renewcommand{\theAppCount}{\Asbuk{AppCount}}

% вставить перед первым \appsection в документ:

%\makeatletter
%\renewcommand{\@currentlabel}{\theAppCount}
%\makeatother

% оформление библиографии и подрисуночных записей через точку
\makeatletter
\renewcommand*{\@biblabel}[1]{\hfill#1.}
\renewcommand*\l@section{\@dottedtocline{1}{1em}{1em}}
\def\redeflsection{\def\l@section{\@dottedtocline{1}{0em}{10em}}}
\makeatother

% библиография
\usepackage[
    backend=biber,
    style=gost-numeric,
    sorting=none,
    movenames=false,
    maxbibnames=99,
    minbibnames=99
]{biblatex}

% Отключаем курсив для авторов в начале записи
\renewcommand*{\mkgostheading}[1]{#1}

% Убираем тире между блоками (в том числе перед городом)
\renewcommand*{\newblockpunct}{\addperiod\space}

% Принудительно задаем маленькую "с" для любых вариантов страниц
\DefineBibliographyStrings{russian}{
	page = {с\adddot},
	pages = {с\adddot},
	pagetotal = {с\adddot}
}

% Настройка оформления списка литературы
\DeclareFieldFormat{labelnumberwidth}{#1.}
\setlength{\biblabelsep}{0.5em}
\defbibenvironment{bibliography}
{\list
	{\printtext[labelnumberwidth]{%
			\printfield{labelprefix}%
			\printfield{labelnumber}}}
	{\setlength{\labelwidth}{0pt}%
		\setlength{\leftmargin}{0pt}%
		\setlength{\labelsep}{0.5em}%
		\setlength{\itemindent}{12.5mm}%
		\addtolength{\itemindent}{\labelsep}%
		\setlength{\itemsep}{\bibitemsep}%
		\setlength{\parsep}{\bibparsep}}%
	\renewcommand*{\makelabel}[1]{\hss##1}}
{\endlist}
{\item}


% Заголовок раздела литературы
\defbibheading{mybib}{%
	\phantomsection
	\anonsection{Список использованных источников}
}
\renewcommand\multicitedelim{\addcomma\space}

% листинги
\usepackage[newfloat,cachedir=minted_cache]{minted}

% настройка листингов
\setminted{
    breaklines=true,
    breakanywhere=true,
    linenos=true,
    numbersep=5pt,
    frame=none,
    xleftmargin=0.25cm
}

% подписи листингов
\captionsetup[listing]{%
	position=above,
	labelsep=dash,
	format=plain,
	justification=raggedright,
	singlelinecheck=false,
	name=Листинг
}

% bash
\setminted[bash]{
    linenos=false,
}

% окружение для длинных листингов
\usepackage{mdframed}
\surroundwithmdframed[
    hidealllines=true,
    innerleftmargin=0pt,
    innerrightmargin=0pt,
    innertopmargin=0pt,
    innerbottommargin=0pt
]{minted}

\usepackage{lastpage}
